\chapter{Intégrale de Lebesgue}
\lecture{29 octobre}
\section{Intégrale de Lebesgue pour des fonctions simples} 
\begin{definition}
	Soit $\varphi$ une fonction simple de représentation canonique $\varphi = \sum_{k = 1}^n a_k \chi_{E_k}$. Alors on pose
	\[
	\int_{\Rn} \varphi(x) \, \mathrm dx = \sum_{k=1}^n a_k m(E_k)
	\]
\end{definition}

\begin{remarque} \label{rem: Invariance par représentation}
	On peut montrer que cette définition est indépendante de la représentation choisie, ainsi si on choisit 
	\[
	\varphi = \sum_{k=1}^m b_k \chi_{\tilde E_k}
	\]
	une représentation arbitraire de $\varphi$, avec $\tilde E_k$ mesurables de mesure finie mais pas nécessairement disjoints, et les $b_k$ par nécessairement distincts. Alors on a quand même
	\[
	\int_{\Rn} \varphi(x) \, \mathrm dx = \sum_{k=1}^m b_k \chi_{\tilde E_k}
	\] 
\end{remarque}

\begin{proposition}[Propriétés de l'intégrale de fonctions simples] \label{prop: Propriétés de l'intégrale de fonctions simples}
	On a les propriétés suivantes
	\begin{enumerate}
		\item (Linéarité) Si $\varphi$ et $\phi$ sont deux fonctions simples sur $\Rn$, alors l'on a
		\[
		\int_{\Rn} (\alpha \varphi + \beta \phi)(x)\, \mathrm dx = \alpha \int_{\Rn} \varphi(x) \, \mathrm dx + \beta \int_{\Rn} \phi(x) \, \mathrm dx
		\]
		pour tout $\alpha, \beta \in \R$.
		\item (Localisation et additivité). Soient $E$ et $F$ des sous-ensembles mesurables de $\Rn$. Si $\varphi$ est simples, alors
		\[
		\int_E \varphi(x) \, \mathrm dx = \int_{\Rn} \varphi(x) \chi_E(x) \, \mathrm dx
		\]
		est bien défini parce que que $\chi_E$ est simple. Alors on a
		\[
		\int_{E \cup F} \varphi(x) \, \mathrm dx = \int_E \varphi(x) \, \mathrm dx + \int_F \varphi(x) \, \mathrm dx
		\]
		pourvu que $E$ et $F$ soient disjoints.
		
		\item (Monotonicité) Si $\varphi$ et $\psi$ sont des fonctions simples avec $\varphi \leqslant \psi$, alors
		\[
		\int_{\Rn} \varphi(x) \, \mathrm dx \leqslant \int_{\Rn} \psi(x) \, \mathrm dx
		\]
		\item (Inégalité triangulaire9 On a
		\[
		\left|\int_{\Rn}\varphi(x) \, \mathrm dx\right| \leqslant \int_{\Rn} |\varphi(x)| \, \mathrm dx
		\]
		pourvu que $\varphi$ soit simple.
	\end{enumerate}
\end{proposition}

\prf{Toutes ces propriétés découlent de la \autoref{rem: Invariance par représentation}, pour le point $1.$, on observe que si $\varphi = \sum_{k=1}^n a_k \chi_{E_k}$ et $\phi = \sum_{l=1}^m b_l \chi_{F_l}$, alors
	\[
	\alpha \varphi + \beta \phi = \sum_{k=1}^m \alpha a_k \chi_{E_k} + \sum_{l=1}^m \beta b_l \chi_{F_l}
	\]
	et donc par la \autoref{rem: Invariance par représentation} 
	\[
	\int_{\Rn} (\alpha \varphi + \beta \phi)(x) \, \mathrm dx = \sum_{k=1}^n \alpha a_k m(E_k) + \sum_{l=1}^m \beta b_l m(F_l) = \alpha \int_{\Rn} \varphi(x) \, \mathrm dx + \beta \int_{\Rn} \phi(x) \, \mathrm dx
	\]
	De manière similaire, on remarque que
	\[
	\chi_E \varphi = \sum_k a_k \chi_{E_k \cap E} \et \chi_F \varphi = \sum_k a_k \chi_{E_k \cap F}
	\]
	et donc par la \autoref{rem: Invariance par représentation} on a
	\begin{align*}
		\int_E \varphi(x) \, \mathrm dx + \int_F \varphi(x) \, \mathrm dx &= \sum_k m(E_k \cap E) + \sum_k a_k m(E_k \cap F) \\
																			&= \sum_k a_k m(E_k \cap (E \cup F)) = \int_{E \cup F} \varphi(x) \, \mathrm dx
	\end{align*}
	où on utilise que $E$ et $F$ sont disjoints par hypothèse. Et pour la $3.$, on remarque que $\psi - \varphi$ est une fonction simple positive, et donc sa représentation canonique a la forme
	\[
	\psi - \varphi = \sum_k c_k \chi_{C_k}, \quad c_k \geqslant 0
	\]
	Il suit que
	\[
	\int_{\Rn} (\psi - \varphi) \, \mathrm dx = \sum_k c_k \chi_{C_k} \geqslant 0
	\]
	et donc par le point 1. on a bien $\int_{\Rn} \psi(x) \, \mathrm dx \geqslant \int_{\Rn} \varphi(x) \, \mathrm dx$.
}
\lecture{2 novembre}
\section{Intégrale de Lebesgue pour des limites de fonctions simples}



\begin{proposition}[Egorov] \label{prop: Egorov}
	Soit $f : \Rn \to \overline \R$ une fonction mesurable prenant des valeurs finies, presque partout, de support $E \subset \Rn$ de mesure finies. Soit $\{f_k\}_{k=1}^\infty$ une suite de fonctions mesurables, tous de support en $E$ tel que $f_k \underset{k \to \infty}{\longrightarrow} f$ presque partout. Alors pour tout $\varepsilon > 0$, il existe $A_\varepsilon \subset E$ tel que
	\[
	m(E \backslash A_\varepsilon) < \varepsilon
	\] 
	et tel que 
	\[
	{f_k}_{\mid A_\varepsilon}
	\]
	converge uniformément vers $f_{\mid A_\varepsilon}$.
\end{proposition}



\prf{Si on enlève un ensemble de mesure nulle de $E$, on peut alors supposer que $f$ prend des valeurs finies partout et que $f_k \ktoinfty f$ partout. Et donc par abus de language, l'ensemble modifié est encore appelé $E$. Donc d'après l'hypothèse,
\[
\forall x \in E, \quad f_k(x) \ktoinfty f(x) \iff \forall x \in E, \, \forall \varepsilon > 0, \, \exists N, \quad |f_k(x) - f(x) | < \varepsilon, \quad \forall k \geqslant N
\]
On va spécialiser $\varepsilon = \frac1k$. Et donc on peut écrire
\[
E = \bigcup_{N \geqslant 1} \{ x \in E \mid |f_n(x) - f(x)| < \frac1k, \quad \forall n  \geqslant N\}
\]
A droite, un a une réunion croissante dénombrable d'ensemble mesurables, et donc par la \autoref{prop: limite de mesurabilité} on a 
\[
m(E) = \lim_{N \to \infty} m(\{ x \in E \mid |f_n(x) - f(x)| < \frac1k, \quad \forall n  \geqslant N\})
\]
Il suit qu'il existe $N_k \in \N$ tel que si on pose $E_k = \{ x \in E \mid |f_n(x) - f(x)| < \frac1k, \quad \forall n  \geqslant N\}$ l'on ait
\[
m(E_k) > m(E) - \frac{\varepsilon}{2^k}
\]
Pour finir, on pose $A_\varepsilon = \bigcap_{k \geqslant 1} E_k$, alors on a 
\[
m(E\backslash A_\varepsilon) \leqslant \sum_{k \geqslant 1} m(E \backslash E_k) \leqslant \sum_{k \geqslant 1} (m(E) - m(E_k)) < \sum_{k \geqslant 1} \frac{\varepsilon}{2^k} = \varepsilon
\]
De plus, soit $\delta > 0$. On choisit $k \geqslant 1$ tel que $\frac1k < \delta$. Par construction $A_\varepsilon \subset E_k$, il suit, par définition, que si $n \geqslant N$, on a $|f_n(x) - f(x)| < \frac1k < \delta$.
}

\begin{theoreme} \label{thm: Limite de suite d'intégrale de fonctions simples}
	Soit $\{\varphi_k\}_{k=1}^\infty$ une suite de fonctions simples, bornée uniformément, et tous de support contenu dans $E \subset \Rn$ avec $m(E) < \infty$. Supposons que $\varphi_k \ktoinfty f$ presque partout. Alors
	\[
	\lim_{k \to \infty} \int_{\Rn} \varphi_k(x) \, \mathrm dx
	\]
	existe. De plus si $f = 0$, on a $\lim_{k \to \infty} \int_{\Rn} \varphi_k(x) \, \mathrm dx = 0$.
\end{theoreme}

\prf{Par la \autoref{prop: Egorov} avec $f = \lim_{k \to \infty} \varphi_k$: on sait que $f$ est mesurables, $f$ est borné presque partout et de support de mesure fini. Alors on peut appliquer la \autoref{prop: Egorov}.
	
Ainsi étant donné $\varepsilon > 0$ arbitraire, il existe $A_\varphi \subset E$, avec $m(E \backslash) < \varepsilon$, et tel que ${\varphi_k}_{\mid A_\varepsilon} \ktoinfty f_{\mid A_\varepsilon}$ uniformément. En particulier, il existe $N \in \N$ tel que
\[
\left| {\varphi_k}_{\mid A_\varepsilon} - {\varphi_{k'}}_{\mid A_\varepsilon}\right| < \varepsilon, \quad \forall k,k' \geqslant N
\]
Et donc de suite
\[
\left|\int_{\Rn} \varphi_k - \varphi_{k'} \, \mathrm dx\right| \leqslant \int_{\Rn} | \varphi_k -\varphi_{k'}| \, \mathrm dx = \int_{A_\varepsilon} | \varphi_k -\varphi_{k'}| \, \mathrm dx + \int_{E \backslash A_{\varepsilon}} | \varphi_k -\varphi_{k'}| \, \mathrm dx \leqslant \varepsilon m(A_\varepsilon) + 2M m(E\backslash A_\varepsilon) \leqslant \varepsilon(m(E) + 2M)
\]
Comme $\varepsilon > 0$ était arbitraire, la convergence de $\{\int_{\Rn} \varphi_k \, \mathrm dx\}$ suit.

Pour la deuxième partie du théorème, on pose $f = 0$, on trouve donc que $|{\varphi_k}_{\mid A_\varepsilon}| < \varepsilon$ et donc
\[
\left| \int_{\Rn} \varphi_k \, \mathrm dx\right| \leqslant \int_{\Rn} |\varphi_k| \, \mathrm dx \leqslant  \int_{A_\varepsilon} |\varphi_k| \, \mathrm dx + M m(E \backslash A_\varepsilon) \leqslant \varepsilon (m(E) + M) 
\]
et donc par arbitrarité de $\varepsilon$ on conclut.
}

\begin{corollaire}
	Soit $f : \Rn \to \R$ une fonction mesurable, bornée, de support de mesure finie. Soit $\{\varphi_k\}_{k = 1}^\infty$ une suite de fonctions simples croissante telle que $\varphi_k \ktoinfty f$ simplement. Alors
	\[
	\lim_{k \to \infty} \int_{\Rn} \varphi_k \, \mathrm dx
	\]
	existe, et est indépendante de la suite $\{\varphi_k\}_{k=1}^\infty$ choisie. On définit donc
	\[
	\int_{\Rn} f \, \mathrm dx := \lim_{k \to \infty} \int_{\Rn} \varphi_k \, \mathrm dx
	\]
\end{corollaire}


\begin{proposition}
	Soit $f : \Rn \to \Rn$ borné, de support compacte, et intégrable au sens de Riemann, alors $f$ est mesurable et sont intégrale au sens de Riemann coïncide avec l'intégrale de Lebesgue.
\end{proposition}


\begin{proposition} \label{prop: Propriétés de l'intégrale de Lebesgue pour des fonctions bornées}
	On a $5$ propriétés de l'intégrale de Lebesgue ; 
	\begin{enumerate}
		\item (Linéarité) Si $f$ et $g$ sont deux fonctions mesurables, bornés et de support de mesure bornée sur $\Rn$, alors on a 
		\[
		\int_{\Rn} (\alpha f + \beta g) \, \mathrm dx = \alpha \int_{\Rn} f(x) \, \mathrm dx + \beta \int_{\Rn} g(x) \, \mathrm dx 
		\]
		\item (Additivité) Soient $E$ et $F$ des sous-ensembles mesurables de $\Rn$. Si $f$ est comme dans le point $1.$, alors
		\[
		\int_{E \cup F} f(x) \, \mathrm dx = \int_E f(x) \, \mathrm dx + \int_F f(x) \, \mathrm dx
		\]
		si on pose $\int_E f(x) \, \mathrm dx = \int_{\Rn} f(x) \chi_E(x) \, \mathrm dx$.
		
		\item (Monotonicité) Si $f,g$ sont deux fonctions comme avant avec $f \leqslant g$, alors
		\[
		\int_{\Rn} f(x) \, \mathrm dx \leqslant \int_{\Rn} g(x) \, \mathrm dx
		\]
		\item (Inégalité triangulaire) On a
		\[
		\left|\int_{\Rn} f(x) \, \mathrm dx \right| \leqslant \int_{\Rn} |f(x)| \, \mathrm dx 
		\]
		\item Si $f = 0$ presque partout, alors
		\[
		\int_{\Rn} f(x) \, \mathrm dx = 0
		\]
		En plus si $f = g$ presque partout, alors
		\[
		\int_{\Rn} f(x) \, \mathrm dx = \int_{\Rn} g(x) \, \mathrm dx 
		\]
	\end{enumerate}
	
\end{proposition}

\prf{Pour le premier point, par le \autoref{thm: suite croissante de fonctions simples} il existe deux suite de fonctions simples croissantes $\{\varphi_k\}_{k \geqslant 1}$ et $\{\psi_k\}_{k \geqslant 1}$, convergeant respectivement vers $f$ et $g$. Alors il est clair que $\alpha \varphi_k + \beta \psi_k$ converge vers $\alpha f + \beta g$ et est borné, et de support de mesure finie uniformément en $k$. Ainsi si $\{\zeta_k\}_{k \geqslant 1}$ est une suite croissante qui converge vers $\alpha f + \beta g$, alors par le \autoref{thm: Limite de suite d'intégrale de fonctions simples} on obtient que
	\[
	\lim_{k \to \infty} \int_{\Rn} (\alpha \varphi_k + \beta \psi_k - \zeta_k) \, \mathrm dx = 0
	\]
	Or par définition, $\int_{\Rn} (\alpha f + \beta g) \, \mathrm dx = \lim_{k \to \infty} \int_{\Rn} \zeta \, \mathrm dx$, et d'après la \autoref{prop: Propriétés de l'intégrale de fonctions simples} on a
	\begin{align*}
		\lim_{k \to \infty} \int_{\Rn} (\alpha \varphi_k + \beta \psi_k) \, \mathrm dx &= \alpha \lim_{k \to \infty} \int_{\Rn} \varphi_k(x) \, \mathrm dx + \beta \lim_{k \to \infty} \int_{\Rn} \psi_k(x) \, \mathrm dx \\ &= \alpha \int_{\Rn} f(x) \, \mathrm dx + \beta \int_{\Rn} g(x) \, \mathrm dx
	\end{align*}
	
Les propriétés $2-4$ sont analogue aux fonctions simples. Donc on regarde juste la dernière affirmation ; si $\varphi_k \to f$ est une suite croissante de fonctions simples, et donc $|\varphi_k| \leqslant |f|$ pour tout $k \geqslant 1$, on a que
\[
\int_{\Rn} \varphi_k(x) \, \mathrm dx = 0, \quad \forall k \geqslant 1
\]
et donc ça entraîne bien
\[
\int_{\Rn} f(x) \, \mathrm dx = 0
\]
}

\lecture{9 novembre}
\begin{theoreme}[Convergence bornée] \label{thm: Convergece bornée}
	Soit $\{f_k\}_{k \geqslant 1}$ une suite de fonctions mesurables et uniformément bornées, et tous de support dans $E \subset \Rn$, où $m(E) < \infty$. Supposons qe $f_k \ktoinfty f$ p.p où $f$ est bornée et de support de mesure finie. Alors $f$ est mesurable et 
	\[
	\lim_{k \to \infty} \int_{\Rn} f_k  \, \mathrm dx = \int_{\Rn} f \, \mathrm dx
	\]	
\end{theoreme}

\prf{On utilise \nameref{prop: Egorov} pour réduire à la convergence uniforme. Soit $\varepsilon > 0$ donné, par la \autoref{prop: Egorov} on choisit $A_\varepsilon \subset E$ tel que
	\[
	{f_k}_{\mid_{A_\varepsilon}} \overset{unif}{\ktoinfty} {f}_{\mid_{A_\varepsilon}}
	\]
	et $m(E \backslash A_\varepsilon) < \varepsilon$. Soit $N_\varepsilon \in \N$ tel que
	\[
	\forall k \geqslant N_\varepsilon, \quad \left|{f_k}_{\mid_{A_\varepsilon}} - {f}_{\mid_{A_\varepsilon}}\right| < \varepsilon
	\]
	On conclut alors que
	\[
	\int_{\Rn} \left| f_k -f \right| \, \mathrm dx \leqslant 	\int_{A_\varepsilon} \left| f_k -f \right| \, \mathrm dx + \int_{E \backslash A_{\varepsilon}} \left| f_k - f\right| \, \mathrm dx \leqslant \varepsilon(m(E) + M + \tilde M)
	\]
	Et donc par l'arbitrarité de $\varepsilon > 0$, la conclusion suit.
}

\section{Intégrale de Lebesgue pour les fonctions générales}

\subsection{Intégrale de Lebesgue pour des fonctions non-négatives}
\begin{definition}
	Soit $f : \Rn \to [0, \infty]$ une fonction mesurable et non-négative. On pose
	\[
	\int_{\Rn} f(x)  \, \mathrm dx = \sup_{0 \leqslant g \leqslant f}\int_{\Rn} g(x) \, \mathrm dx \in \overline \R_+
	\]
	où le suprémum est sur toutes les fonctions $0 \leqslant g \leqslant f$ bornée et mesurables, et de support avec mesure finie. Si on dira que $f$ est \textbf{intégrale au sens de Lebesgue} pour vu que
	\[
	\int_{\Rn} f \, \mathrm dx < \infty
	\]
\end{definition}

\begin{proposition} \label{prop: propriétés de l'intégrale de Lebesgue pour des fonctions non-négative}
	On a les propriétés suivantes ; 
	\begin{enumerate}
		\item Si $a, b \geqslant 0$ et $f,g : \Rn \to [0, \infty]$ sont mesurables. Alors
		\[
		\int_{\Rn} a f + b g \, \mathrm dx = a \int_{\Rn} f \, \mathrm dx + b \int_{\Rn} g \, \mathrm dx
		\]
		\item  Si $E, F \subset \Rn$ sont mesurables, disjoints, et que l'on pose $\int_E f \, \mathrm dx = \int_{\Rn} \chi_X f \, \mathrm dx$, alors on a
		\[
		\int_E f \, \mathrm dx + \int_F f \, \mathrm dx = \int_{E \cup F} f \, \mathrm dx.
		\]
		\item Si $0 \leqslant f_1 \leqslant f_2$ sont mesurables, alors
		\[
		\int_{\Rn} f_1 \, \mathrm dx \leqslant \int_{\Rn} f_2 \, \mathrm dx
		\]
		\item Si $f= 0$ p.p, alors
		\[
		\int_{\Rn} f \, \mathrm dx = 0
		\]
		\item Si $f \geqslant 0$ est intégrale, on a que $f(x) \in [0,\infty[$ p.p
	\end{enumerate}
\end{proposition}

\prf{Dans l'ordre on a
	\begin{enumerate}
		\item Pour commencer, soit $a \geqslant 0$, alors
		\begin{align*}
			\int_{\Rn} a f(x) \, \mathrm dx &= \sup_{0 \leqslant g \leqslant af}\int_{\Rn} g(x) \, \mathrm dx = \sup_{0 \leqslant ag \leqslant af}\int_{\Rn} ag(x) \, \mathrm dx \\ &= a \sup_{0 \leqslant g \leqslant f}\int_{\Rn} g(x) \, \mathrm dx = a \int_{\Rn} f(x) \, \mathrm dx
		\end{align*}
		Il reste donc à montrer l'additivité ; soit $f,g : \Rn \to [0, \infty]$ mesurables. Or par définition
		\[
		\int_{\Rn} f(x) \, \mathrm dx = \sup_{0 \leqslant \varphi \leqslant f} \int_{\Rn} \varphi(x) \, \mathrm dx \et \int_{\Rn} g(x) \, \mathrm dx = \sup_{0 \leqslant \psi \leqslant g} \int_{\Rn} \psi(x) \, \mathrm dx
		\]
		Ainsi soit $0 \leqslant \varphi \leqslant f$ et $0 \leqslant \psi \leqslant g$, alors $0 \leqslant \varphi + \psi \leqslant f + g$ et donc
		\[
		\int_{\Rn} \varphi(x) \, \mathrm dx + \int_{\Rn} \psi(x) \, \mathrm dx \leqslant \int_{\Rn} (f+g) \, \mathrm dx
		\]
		et donc par passage au suprémum
		\[
		\int_{\Rn} f(x) \, \mathrm dx + \int_{\Rn} g(x) \, \mathrm dx \leqslant \int_{\Rn} (f+g)(x) \, \mathrm dx
		\]
		Inversement, soit $0 \leqslant \zeta \leqslant f+ g$. On construit $\varphi$ et $\psi$ tel que $\varphi + \psi = \zeta$ et, $0\leqslant \varphi \leqslant f$ et $0 \leqslant \psi \leqslant g$. Il alors que
		\[
		\int_{\Rn}\zeta(x) \, \mathrm dx = \int_{\Rn} \varphi(x) \, \mathrm dx + \int_{\Rn} \psi(x) \, \mathrm dx \leqslant \int_{\Rn} f(x) \, \mathrm dx + \int_{\Rn} g(x) \, \mathrm dx
		\]
		et donc par passage au suprémum,
		\[
		\int_{\Rn} (f+g)(x) \, \mathrm dx \leqslant \int_{\Rn} f(x) \, \mathrm dx + \int_{\Rn} g(x) \, \mathrm dx
		\]
		Et donc par double inégalité le propriété suit.
		
		\item[5.] Supposons que $f : \Rn \to [0, \infty]$ soit intégrale. On note $E := \{x \in \Rn \mid f(x) = \infty\}$ et supposons que $m(E) > 0$. En posant $E_k = E \cap B(0, k)$ alors par convergence de $E_k$ on peut choisir $k_* \geqslant 1$ tel que $0 < m(E_{k_+}) < \infty$. On poser donc
		\[
		\varphi_k := k \chi_{E_{k_+}}, \quad k \geqslant 1
		\]
		et donc
		\[
		\int_{\Rn} f(x) \, \mathrm dx \geqslant \int_{\Rn} \varphi_k(x) \, \mathrm dx = km(E) \ktoinfty \infty
		\]
	\end{enumerate}
}

\subsection{Théorème de convergence pour l'intégrale de Riemann au cadre des fonctions non-négative}

\begin{theoreme}[Lemme de Fatou] \label{thm: Lemme de Fatou}
	Soit $\{f_k\}_{k \geqslant 1}$ une suite de fonctions non-négative et mesurables sur $\Rn$ tel que $f_k \ktoinfty f$, alors on a
	\[
	\int_{\Rn} f(x) \, \mathrm dx \leqslant \liminf_{k \to \infty} \int_{\Rn} f_k(x) \, \mathrm dx 
	\]
\end{theoreme}

\prf{Soit $g$ une fonction non-négative, bornée, de support de mesure fini, et $0 \leqslant g \leqslant f$. Alors on pose
	\[
	\varphi_k = := \min\{g, f_k\}
	\]
	qui sont bornées, de support de mesure finie, et avec $\varphi_k \leqslant f_k$. Il suit alors que
	\[
	\int_{\Rn} \varphi_k(x) \, \mathrm dx \leqslant \int_{\Rn} f_k(x) \, \mathrm dx 
	\]
	Or par hypothèse de convergence $\varphi_k \ktoinfty g$ parce que $g \leqslant f$, et les fonctions $\varphi_k$ sont bornées. Ainsi par le \autoref{thm: Convergece bornée} on a
	\[
	\lim_{k \to \infty} \int_{\Rn} \varphi_k(x) \, \mathrm dx = \int_{\Rn} g(x) \, \mathrm dx
	\]
	Et donc 
	\[
	\int_{\Rn} g(x) \, \mathrm dx = \lim_{k \to \infty} \int_{\Rn} \varphi_k(x) \, \mathrm dx = \liminf_{k \to \infty} \int_{\Rn} \varphi_k(x) \, \mathrm dx \leqslant \liminf_{k \to \infty} \int_{\Rn} f_k(x) \, \mathrm dx 
	\]
	Finalement, par passage au suprémum le théorème suit.
}
\begin{lemme} \label{lem: Convergence d'une suite bornée par f}
	Soit $0 \leqslant f_k \leqslant f$, une suite de fonction mesurables sur $\Rn$ tel que $f_k \ktoinfty f$, alors on a
	\[
	\lim_{k \to \infty} \int_{\Rn} f_k(x) \, \mathrm dx = \int_{\Rn} f(x) \, \mathrm dx
	\]
\end{lemme}

\prf{Par le \nameref{thm: Lemme de Fatou} et la monotonie de l'intégrale de Lebesgue on a 
	\[
	\int_{\Rn} f(x) \, \mathrm dx \leqslant \liminf_{k \to \infty} \int_{\Rn} f_k(x) \, \mathrm dx \et \limsup_{k \to \infty} \int_{\Rn} f_k(x) \, \mathrm dx \leqslant \int_{\Rn} f(x) \, \mathrm dx
	\]
	Donc le lemme suit.
}
\begin{theoreme}[Convergence monotone] \label{thm: Convergence monotone}
	Soit $\{f_k\}_{k \geqslant 1}$ une suite croissante de fonctions mesurables et non-négative	telle que $f_k \ktoinfty f$, alors on a
	\[
	\lim_{k \to \infty} \int_{\Rn} f_k(x) \, \mathrm dx = \int_{\Rn} f(x) \, \mathrm dx
	\]
\end{theoreme}

\prf{Cas particulier du \autoref{lem: Convergence d'une suite bornée par f}.
}

\lecture{12 novembre}

\subsection{Intégrale de Lebesgue au cas général}

\begin{definition}
	Soit $f : \Rn \to [-\infty, \infty]$ une fonction mesurable. Alors $f$ est dit \textbf{intégrable au sens de Lebesgue}, pourvu que $|f|$ est intégrable. Et on pose
	\[
	\int_{\Rn} f(x) \, \mathrm dx := \int_{\Rn} f_+(x) \, \mathrm dx - \int_{\Rn} f_-(x) \, \mathrm dx 
	\]
\end{definition}

\begin{lemme}
	Si $f = f_1 - f_2$ où les fonctions $f_{1,2}$ sont mesurables, non-négatives et intégrables, alors $f$ est intégrable, et on a
	\[
	\int_{\Rn} f(x) \, \mathrm dx = \int_{\Rn} f_1(x) \, \mathrm dx - \int_{\Rn} f_2(x) \, \mathrm dx 
	\]
\end{lemme}

\prf{On commence par remarquer que $f$ est intégrable puisque $|f| \leqslant f_1 + f_2$. De suite on a que
	\[
	f = f_+ - f_- = f_1 - f_2
	\]
	on en déduit que $f_+ + f_2 = f_- + f_1$, ainsi par la \autoref{prop: propriétés de l'intégrale de Lebesgue pour des fonctions non-négative} on a
	\[
	\int_{\Rn} f_+ \, \mathrm dx + \int_{\Rn} f_2 \, \mathrm dx = \int_{\Rn} f_- \, \mathrm dx + \int_{\Rn} f_1 \, \mathrm dx 
	\]
	Ce qui entraine bien que
	\[
	\int_{\Rn} f_1(x) \, \mathrm dx - \int_{\Rn} f_2(x) \, \mathrm dx = \int_{\Rn} f_+(x) \, \mathrm dx - \int_{\Rn} f_-(x) \, \mathrm dx = \int_{\Rn} f(x) \, \mathrm dx	
	\]
}

\begin{proposition}
	Les fonctions intégrables au sens de Lebesgue forment un espace vectoriel sur $\R$. Et les propriétés 1. - 4. de la \autoref{prop: Propriétés de l'intégrale de Lebesgue pour des fonctions bornées} s'appliquent aux intégrales de fonctions générales.
\end{proposition}

\begin{theoreme}[Théorème de convergence dominée] \label{thm: Théorème de convegence dominée}
	Soit $f_k : \Rn \to [-\infty, \infty]$ une suite de fonctions mesurables telle que $f_k \ktoinfty f$. S'il existe une fonction $g : \Rn \to [0, \infty]$ intégrable telle que
	\[
	\forall k \geqslant 1, \quad |f_k| \leqslant g
	\]
	Alors on a la convergence des intégrales
	\[
	\int_{\Rn} f_k (x) \, \mathrm dx \ktoinfty \int_{\Rn} f(x) \, \mathrm dx
	\]
\end{theoreme}

\prf{Soit $\varepsilon > 0$. Pour tout $N \geqslant 1$, on définit
	\[
	E_N := \{ x \in \Rn \mid |x| < N \et g(x) \leqslant N\}
	\]
	Évidement les ensembles $E_N$ forment une suite croissante, et donc
	\[
	\chi_{E_N} g \to g
	\]
	est la limite d'une suite croissante de fonction. Ainsi par le \autoref{thm: Convergence monotone} on a que
	\[
	\lim_{N \to \infty} \int_{\Rn} \chi_{E_N} g \, \mathrm dx = \int_{\Rn} g \, \mathrm dx < \infty
	\]
	on en déduit que
	\[
	\lim_{N \to \infty} \int_{\Rn} \chi_{E_N^c} g \, \mathrm dx = 0
	\]
	et donc il existe $N_\varepsilon \in \N$ tel que
	\[
	\int_{\Rn}  \chi_{E_N^c} g \, \mathrm dx < \frac\varepsilon3, \quad \forall N \geqslant N_\varepsilon
	\]
	Et donc par hypothèse
	\[
	\int_{\Rn} \chi_{E_N^c} |f_k| \, \mathrm dx < \frac\varepsilon3 \et \int_{\Rn} \chi_{E_N^c} |f| \, \mathrm dx < \frac\varepsilon3
	\]
	où $N \geqslant N_\varepsilon$. Or le suite $\{\chi_{E_{N_\varepsilon}} f_k\}_{k \geqslant 1}$ est uniformément bornée et de support de mesure finie en $k \geqslant 1$, et donc par le \autoref{thm: Convergece bornée}, on a
	\[
	\lim_{k \to \infty} \int_{\Rn} \chi_{E_{N_\varepsilon}} |f_k - f | \, \mathrm dx = 0
	\]
	Il en suite que pour tout $k \geqslant K_\varepsilon$, on a
	\[
	\int_{\Rn} \chi_{\E_{N_\varepsilon}} |f_k - f |\, \mathrm dx < \frac\varepsilon3
	\]
	Ce qui nous donne finalement, pourvu que $k \geqslant K_\varepsilon$,
	\begin{align*}
	\int_{\Rn} |f_k - f| \, \mathrm dx &= \int_{\Rn} \chi_{E_{N_\varepsilon}} |f_k - f| \, \mathrm dx + \int_{\Rn} \chi_{E_{N_\varepsilon}^c} |f_k - f| \, \mathrm dx \\
										& \leqslant \int_{\Rn} \chi_{E_{N_\varepsilon}} |f_k - f| \, \mathrm dx + \int_{\Rn} \chi_{E_{N_\varepsilon}^c} |f_k| \, \mathrm dx + \int_{\Rn} \chi_{E_{N_\varepsilon}^c} |f| \, \mathrm dx \\
										& < \varepsilon
	\end{align*}
	et donc le théorème suit.
}

\section{l'espace $\L1$}

\begin{definition}
	On définit l'espace $\L1$ comme
	\[
	\L1 := \big\{ f : \Rn \to \C ~\big|~ \int_{\Rn} |f(x)| \,  \mathrm dx < \infty \big\}
	\]
\end{definition}

\begin{proposition}
	Soit $\L1$ la collection de classe d'équivalence de fonctions intégrales sur $\Rn$, où $f \sim g$ pourvu que 
	\[
	f (x) = g(x)
	\]
	presque partout. Alors $\L1$ est un espace vectoriel normé muni de la norme $\nl1{\cdot}$.
\end{proposition}

\prf{On remarque que si $f \sim f'$ et $g \sim g'$ alors
	\[
	f + g \sim f' + g'
	\]
	Ainsi si $[f]$ note la classe d'équivalence de $f$, on peut donc définir $[f] + [g] := [f + g]$, parce que le coté droit ne dépend pas de représentant $f,g$. En plus on peut poser
	\[
	\nl1{[f]} := \nl1f
	\]
	parce que le côté droit ne change pas si $f$ est remplacé par un autre représentant de la même classe d'équivalence. Et finalement on a bien
	\[
	\nl1{f + g} = \int_{\Rn} |f+g| \dx \leqslant \int_{\Rn} |f| \dx + \int_{\Rn} |g| \dx = \nl1f + \nl1g
	\]
	ainsi que
	\[
	\nl1f = \int_{\Rn} |f| \dx = 0
	\]
	si et seulement si $f = 0$ p.p. D'où la proposition.
}


\begin{theoreme}[Complétude]
	L'espace $\L1$, muni de la norme $\nl1\cdot$ est complet, et donc forme un espace de Banach.
\end{theoreme}

\prf{Soit $\{f_k\}_{k \geqslant 1}$ une suite de Cauchy dans $\L1$. Il faut donc montrer qu'il existe $f \in \L1$ telle que
	\[
	\lim_{k \to \infty} \nl1{f_k-f} = 0
	\]
	Comme la suite des $f_k$ est de Cauchy, on peut extraire une suite-suite $\{f_{n_k}\}_{k \geqslant1}$ telle que
	\[
	\forall k \geqslant 1, \quad \nl1{f_{n_{k+1}} - f_{n_k}} < 2^{-k}
	\]
	On pose
	\[
	g_j := |f_{n_1}| + \sum_{l=1}^j |f_{n_{j+1}} - f_{n_l}|, \quad j \geqslant 1
	\]
	Il est claire que la suite $g_j$ est croissante, non-négative, et mesurable, et donc par le \autoref{thm: Convergence monotone}, on a
	\[
	\lim_{j \to \infty} \int_{\Rn} g_j(x) \dx = \int_{\Rn} g(x) \, \mathrm dx
	\]
	où $g := \lim_{j \to \infty} g_j$. Mais par définit de $g_j$ on a
	\[
	\int_{\Rn} g_j(x) \dx \leqslant \nl1{f_{n_1}} + \sum_{l=1}^j 2^{-l} < \nl1{f_{n_1}} + 1
	\]
	Ainsi la fonction $g = |f_{n_1}| + \sum_{l=1}^\infty |f_{n_{j+1}} - f_{n_l}|$ est intégrable. Alors par la \autoref{prop: propriétés de l'intégrale de Lebesgue pour des fonctions non-négative} on ne déduit que
	\[
	|f_{n_1}(x)| + \sum_{l=1}^\infty |f_{n_{j+1}}(x) - f_{n_l}(x)| < \infty
	\]
	presque partout. Pour un tel $x$ où on a convergence, la série
	\[
	f_{n_1}(x) + \sum_{l=1}^\infty f_{n_{l+1}}(x) - f_{n_l}(x) = \lim_{j \to \infty} \Big( f_{n_1}(x) + \sum_{l=1}^j f_{n_{l+1}}(x) - f_{n_l}(x)\Big) = \lim_{j \to \infty} f_{n_{j+1}}(x)
	\]
	converge aussi. Si on appelle $f$ cette limite, qui est donc mesurable, on a que
	\[
	\lim_{j \to \infty} \nl1{f - f_{n_{j+1}}} = 0
	\]
	par le \autoref{thm: Théorème de convegence dominée}, puisque
	\[
	\forall j \geqslant 1, \quad |f_{n_{j+1}}| \leqslant g
	\]
	et $g$ est intégrable. Finalement soit $\varepsilon > 0$, on choisit $N_\varepsilon \in \N$ assez grand tel que
	\[
	\forall p,q \geqslant N_\varepsilon, \quad \nl1{f_p - f_q} \frac\varepsilon2
	\]
	De plus par convergence de la suite $\{f_{n_j}\}$ vers $f$, il existe $k \geqslant N_\varepsilon$ et $n_{j_\varepsilon} > k$ tel que
	\[
	\nl1{f - f_{n_\varepsilon}} < \frac\varepsilon2
	\]
	Et donc on conclut que
	\[
	\nl1{f - f_k} \leqslant \nl1{f - f_{n_{j_{\varepsilon}}}} +\nl1{f_k -f_{n_{j_{\varepsilon}}}} < \varepsilon
	\]
	Ainsi le théorème suit.
}

\subsection{Sous-espace dense dans $\L1$}

\begin{proposition}
	Les espaces de fonctions suivants sont denses dans $\L1$ : 
	\begin{enumerate}
		\item Les fonctions simples,
		\item Les combinaisons linéaires de fonctions caractéristiques de boîtes rectangulaires,
		\item Les fonctions continues de support compacte, $C_c^0(\Rn)$.
	\end{enumerate}
\end{proposition}

\prf{\begin{enumerate}
		\item On traite que le cas de $f$ réel. De plus en utilisant la représentation $f = f_+ - f_1$, on est réduit au cas $f \geqslant 0$. D'après le \autoref{lem: suite de fonctions mesurables} il existe une suite croissante $\{\varphi_k\}_{k \geqslant 1}$ de fonctions simples non-négatives telle que
		\[
		\varphi \ktoinfty f
		\]
		et donc par le \autoref{thm: Convergence monotone} on en déduit que
		\[
		\nl1f = \int_{\Rn} f(x) \dx = \lim_{k \to \infty} \int_{\Rn} \varphi_k(x) \dx
		\]
		et comme $0 \leqslant \varphi_K \leqslant f$ ceci implique que
		\[
		\lim_{k \to \infty} \nl1{f - \varphi_k} = \lim_{k \to \infty} \int_{\Rn} f - \varphi_k \dx = \lim_{k \to \infty} \left( \int_{\Rn} f(x) \dx - \int_{\Rn} \varphi_k(x) \dx\right) = 0 
		\]
		
		\item D'après le point 1., il suffit de montrer que tout fonction simples peut être approximée par des combinaisons linéaires de fonctions caractéristique de boîtes rectangulaires, et de suite que $\chi_E$ peut être approximé ainsi, où $E$ est mesurable et de mesure finie. Par définition de la mesure \autoref{def: Mesure de Lebesgue}, étant donné $\varepsilon > 0$, il existe une collection de cube $\{Q_j\}_{j \geqslant 1}$ telle que
		\[
		E \subset \bigcup_j Q_j \et \sum_j \Vol Q_j < m(E) + \frac\varepsilon4
		\]
		Soit $N \in \N$ tel que $\sum_{j \geqslant N+1} \Vol Q_j < \frac{\varepsilon}4$. On a alors que
		\[
		m \Bigg(\bigcup_{j=1}^N \backslash E\Bigg) < \frac\varepsilon4 \et m \Bigg( E \backslash \bigcup_{j=1}^N Q_j \Bigg) < \frac\varepsilon4
		\]
		ce qui implique que
		\[
		\nl1{\chi_E - \chi_{\bigcup_{j=1}^N Q_j}} < \frac\varepsilon2
		\]
		Or on peut écrire
		\[
		\bigcup_{j=1}^N Q_j = \bigcup_{l=1}^M R_l
		\]
		pour une collection $\{R_l\}_{l=1}^M$ de boîtes rectangulaires presque disjointes, et de la remplacer par des boîtes rectangulaires un peu plus petites et disjointes $\{ \tilde R_l\}_{l=1}^M$, de telle que manière que
		\[
		\nl1{\chi_{\bigcup_{j=1}^N Q_j} - \sum_{l=1}^M\chi_{\tilde R_l}} < \frac\varepsilon2
		\]
		Il suit que
		\[
		\nl1{\chi_E - \sum_{l=1}^M \chi_{\tilde R_l}} \leqslant \nl1{\chi_E - \chi_{\bigcup_{j=1}^N Q_j}} + \nl1{\chi_{\bigcup_{j=1}^N Q_j} - \sum_{l=1}^M\chi_{\tilde R_l}} < \varepsilon
		\]
		
		\item D'après le point 2., il suffit d'approximer des fonctions caractéristiques de boîtes rectangulaires par des fonctions dans $C_c^0(\Rn)$. En dimension $1$, si $I = [a,b]$, étant donnée $\varepsilon > 0$, on peut approximer la fonction caractéristique $\chi_I$ par la fonction continue et de support compact
		\begin{figure}[H]
		\centering
		\begin{tikzpicture}[scale=1.2]
			
			% Paramètres
			\def\a{1.5}
			\def\b{3.5}
			\def\e{0.8}
			
			% Axes
			\draw[-{Stealth}] (-0.3,0) -- (5.5,0) node[right] {$x$};
			\draw[-{Stealth}] (0,-0.3) -- (0,1.5) node[above] {$\phi(x)$};
			
			% Graphe de phi
			\draw[thick, blue]
			(0,0) -- (\a-\e,0) -- (\a,1) -- (\b,1) -- (\b+\e,0) -- (5.3,0);
			
			% Pointillés verticaux
			\draw[dashed, gray] (\a-\e,0) -- (\a-\e,-0.1);
			\draw[dashed, gray] (\a,1) -- (\a,0);
			\draw[dashed, gray] (\b,1) -- (\b,0);
			
			% Graduations sur l'axe x
			\foreach \x/\lab in {\a-\e/{a-\varepsilon}, \a/a, \b/b, \b+\e/{b+\varepsilon}} {
				\draw (\x,0.05) -- (\x,-0.05);
				\node[below] at (\x,-0.05) {$\lab$};
			}
			
			% Graduation sur l'axe y
			\draw (0.05,1) -- (-0.05,1) node[left] {$1$};
			
		\end{tikzpicture}
		\end{figure}
		Dont on vérifie facilement que
		\[
		\nl1{\chi_I - \phi} = \varepsilon
		\]
		En dimension $n \geqslant 2$, on peut utiliser un produit de fonction comme $\phi$ par rapport à chaque coordonnée pour approximer $\chi_R$.
	\end{enumerate}
}

\subsection{Applications; propriétés d'invariance et de continuité}